\documentclass[11pt]{article}

\usepackage[margin=1in]{geometry}
\usepackage{amsmath, amssymb}
\usepackage{enumitem}
\usepackage{actuarialangle}
\usepackage{setspace}

\setlength{\parindent}{0pt}
\setlength{\parskip}{8pt}

\begin{document}

\begin{center}
    {\Large \textbf{Bonds Practice}}\\
\end{center}

\vspace{10pt}

\hrule

\medskip

\textbf{Problem 1.} Mario has just purchased a \$10{,}000 par value 20-year bond for \$8994.32 with the following characteristics:

\begin{itemize}
    \item The bond has an annual nominal coupon rate of \(r\%\) payable semiannually.
    \item The bond has an annual nominal yield rate of \(9\%\) convertible semiannually.
    \item The bond has a redemption value of \$15{,}000.
\end{itemize}

Calculate the book value of the bond after 12.5 years immediately after the coupon payment.

\medskip
\hrule


\textbf{Problem 2.} You are given the following information on a bond:

\begin{enumerate}[label=\roman*.]
    \item Par value = 1000
    \item Redemption value = 1000
    \item Coupon rate = 12 percent, convertible semiannually
    \item It is priced to yield 10 percent, convertible semiannually
\end{enumerate}

The bond has a term of \(n\) years. If the term of the bond is doubled, the price will increase by 50.

Calculate the price of the \(n\)-year bond.

\medskip
\hrule


\textbf{Problem 4.} Bond A has a par value of \$1{,}000 and pays coupons semiannually at a nominal annual rate of \(8\%\). Bond A has a redemption value of \(C\) and matures in 12 years.

Bond B also has a par value of \$1{,}000, pays coupons semiannually at a nominal annual rate of \(r\), and has the same redemption value \(C\). Bond B also matures in 12 years.

Both bonds have a yield rate of \(6\%\) compounded semiannually.

It is given that the price of Bond B is \$150 less than the price of Bond A.

Calculate \(r\).

\medskip
\hrule


\begin{center}
\begin{large}
\textbf{Challenging}!
\end{large}
\end{center}


\textbf{Problem I.} Kelly purchases a 25-year annual coupon bond for \$20{,}000 with an annual effective yield rate of \(i\). The amount for the accumulation of discount in the first coupon is 50.

The bond's book value is \$20{,}491.51 just after the \(8^\text{th}\) coupon payment.

Calculate the bond's redemption value.

\medskip
\medskip

\begin{enumerate}[label=\Alph*.]
    \item \$22{,}135
    \item \$22{,}170
    \item \$22{,}234
    \item \$22{,}671
    \item \$23{,}660
\end{enumerate}

\medskip
\hrule


\textbf{Problem II.} Three 10-year bonds with par value \$1{,}000 and redemption value equal to par are purchased to yield the same effective rate.

The first bond has a coupon rate of \(8\%\) convertible semiannually and is purchased at a discount of \(X\).

The second bond has a coupon rate of \(9\%\) convertible semiannually and is purchased at a premium of \(Y\).

The third bond has a coupon rate of \(10\%\) convertible semiannually and is purchased at a premium of \(2X\).

Calculate \(Y\).

\begin{enumerate}[label=\Alph*.]
    \item \(\frac{1}{4}X\)
    \item \(\frac{1}{3}X\)
    \item \(\frac{1}{2}X\)
    \item \(\frac{3}{4}X\)
    \item \(X\)
\end{enumerate}

\medskip
\hrule

\textbf{Problem III.} You are given the following information about Bond A and Bond B:

\begin{enumerate}[label=\roman*.]
    \item Bond A is a 20-year bond that has \$10{,}000 face amount, an annual coupon rate of \(7\%\) payable semiannually, and has a redemption value that is \(10\%\) higher than its price. Bond A is bought to yield an annual nominal interest rate of \(5\%\) convertible semiannually.
    
    \item Bond B is an \(n\)-year bond with semiannual coupons and the same face amount, redemption value, coupon rate, and yield rate as Bond A. The price of Bond B is \$158.33 less than the price of Bond A.
\end{enumerate}

Calculate \(n\).

\begin{enumerate}[label=\Alph*.]
    \item 12
    \item 17
    \item 24
    \item 36
    \item 48
\end{enumerate}

\medskip
\hrule

\newpage

\begin{center}
\begin{LARGE}
\textbf{Callable Bonds}
\end{LARGE}
\end{center}


\textbf{Problem I.} Toby purchased a 20-year par value bond with semiannual coupons at a nominal annual rate of \(8\%\) convertible semiannually at a price of \$1{,}722.25. The bond can be called at par value \$1{,}100 on any coupon date starting at the end of year 15.

What is the minimum yield that Toby could receive, expressed as a nominal annual rate of interest convertible semiannually?

\begin{enumerate}[label=\Alph*.]
    \item \(3.2\%\)
    \item \(3.3\%\)
    \item \(3.4\%\)
    \item \(3.5\%\)
    \item \(3.6\%\)
\end{enumerate}

\medskip
\hrule

\textbf{Problem II.} Zach is going to purchase a \$5{,}000-par callable bond that will pay semiannual coupons at an annual rate of \(9\%\). The bond will redeem at par at the end of 20 years.

The bond may be called at the end of year 9 through year 15 with a call price of \$5{,}300.

What is the highest price Zach should pay to yield at least a nominal rate of \(8\%\), convertible semiannually, at redemption?

\begin{enumerate}[label=\Alph*.]
    \item \$5{,}399
    \item \$5{,}434
    \item \$5{,}464
    \item \$5{,}499
    \item \$5{,}524
\end{enumerate}

\medskip
\hrule

\textbf{Problem III.} Matt purchased a 20-year par value bond with semiannual coupons at a nominal annual rate of \(8\%\) convertible semiannually at a price of \$1{,}722.25. The bond can be called at par value \(X\) on any coupon date starting at the end of year 15 after the coupon is paid. The lowest yield rate that Matt can possibly receive is a nominal annual interest rate of \(6\%\) convertible semiannually.

Calculate \(X\).

\begin{enumerate}[label=\Alph*.]
    \item \(1400\)
    \item \(1420\)
    \item \(1440\)
    \item \(1460\)
    \item \(1480\)
\end{enumerate}

\medskip
\hrule

\newpage

\textbf{Problem IV.} You are given the following information about an \(n\)-year bond, where \(n>10\):

\begin{enumerate}[label=\roman*.]
    \item The bond pays \(8\%\) semiannual coupons and has face amount \$1{,}000.
    \item The bond is redeemable at par.
    \item The bond is callable at par 5 years after issue or 10 years after issue.
    \item \(P\) is the price to guarantee a yield of \(6.8\%\) convertible semiannually and \(Q\) is the price to guarantee a yield of \(8.8\%\) convertible semiannually.
    \item \(|P-Q|=123.36\).
\end{enumerate}

Calculate \(n\).

\begin{enumerate}[label=\Alph*.]
    \item \(11\)
    \item \(15\)
    \item \(19\)
    \item \(22\)
    \item \(26\)
\end{enumerate}

\medskip
\hrule

\newpage

\textbf{Problem V.} Bond X is a 20-year bond with annual coupons and the following characteristics:

\begin{enumerate}[label=\roman*.]
    \item Par value is \$1{,}000.
    \item The annual coupon rate is \(10\%\).
    \item Bond X is callable at par on any of the last five coupon dates.
\end{enumerate}

Calculate the maximum purchase price for Bond X that will guarantee an annual effective yield rate of at least \(5\%\).

\begin{enumerate}[label=\Alph*.]
    \item \(1{,}519\)
    \item \(1{,}542\)
    \item \(1{,}570\)
    \item \(1{,}596\)
    \item \(1{,}623\)
\end{enumerate}

\medskip
\hrule

\textbf{Problem VI.} A 1000 par value bond pays annual coupons of 80. The bond is redeemable at par in 30 years, but is callable anytime from the end of the 10th year at 1050.

Based on her desired yield rate, an investor calculates the following potential purchase prices, \(P\):

\begin{itemize}
    \item Assuming the bond is called at the end of the 10th year, \(P = 957\).
    \item Assuming the bond is held until maturity, \(P = 897\).
\end{itemize}

The investor buys the bond at the highest price that guarantees she will receive at least her desired yield rate regardless of when the bond is called.

The investor holds the bond for 20 years, after which time the bond is called.

Calculate the annual yield rate the investor earns.

\begin{enumerate}[label=\Alph*.]
    \item \(8.56\%\)
    \item \(9.00\%\)
    \item \(9.24\%\)
    \item \(9.53\%\)
    \item \(9.99\%\)
\end{enumerate}

\medskip
\hrule

\textbf{Problem VII.} Toby purchased a 20-year par value bond with semiannual coupons of 40 and a redemption value of \$1{,}100. The bond can be called at \$1{,}200 on any coupon date prior to maturity, starting at the end of year 15.

Calculate the maximum price of the bond to guarantee that Toby will earn an annual nominal interest rate of at least \(6\%\) convertible semiannually.

\begin{enumerate}[label=\Alph*.]
    \item \(1{,}251\)
    \item \(1{,}262\)
    \item \(1{,}278\)
    \item \(1{,}286\)
    \item \(1{,}295\)
\end{enumerate}

\medskip
\hrule

\textbf{Problem VIII.} A \$10{,}000 par value bond with \(6\%\) semiannual coupons was sold on 7/1/2005.

It is callable on the following dates:

\begin{center}
\begin{tabular}{|c|c|}
\hline
\textbf{Call Date} & \textbf{Call Price} \\
\hline
7/1/2007 & 11{,}000 \\
1/1/2008 & 11{,}000 \\
7/1/2008 & 11{,}000 \\
1/1/2009 & 10{,}850 \\
7/1/2009 & 10{,}850 \\
1/1/2010 & 10{,}300 \\
7/1/2010 & 10{,}300 \\
1/1/2011 & 10{,}300 \\
7/1/2011 & 10{,}300 \\
1/1/2012 & 10{,}150 \\
7/1/2012 & 10{,}150 \\
1/1/2013 & 10{,}150 \\
\hline
\end{tabular}
\end{center}

If not called, it will be redeemed at par on 7/1/2013.

Calculate the price that ensures a buyer of a nominal yield of \(5\%\) convertible semiannually?

\begin{enumerate}[label=\Alph*.]
    \item \(10{,}638\)
    \item \(11{,}094\)
    \item \(10{,}188\)
    \item \(11{,}122\)
    \item \(10{,}513\)
\end{enumerate}

\medskip
\hrule

\end{document}